\section{Introduction}

Diffusion-weighted magnetic resonance imaging (DWI) is a cornerstone of prostate cancer diagnosis, incorporated into the Prostate Imaging Reporting and Data System (PI-RADS) as a primary determinant of clinical significance \citep{turkbey2019prostate}.
The apparent diffusion coefficient (ADC), derived from monoexponential fitting of the DWI signal at two or more b-values, has become the most widely used quantitative biomarker for prostate cancer detection and grading \citep{manetta2019correlation, yan2024value}.
Lower ADC values indicate restricted water diffusion associated with increased cellular density in malignant tissue, and ADC correlates inversely with Gleason grade---though the lesion-mean value separates adjacent grade groups only modestly \citep{donati2014prostate, rosenkrantz2015whole}.

The monoexponential model, however, collapses the multi-b signal decay into a single scalar, discarding information about the heterogeneous tissue microstructure that generates the observed signal \citep{fennessy2023quantitative}.
Extended b-value protocols spanning $b = 0$ to $3500\;\text{s/mm}^2$ reveal non-monoexponential signal behavior arising from multiple diffusion compartments---including restricted intracellular water, hindered extracellular diffusion, glandular luminal water, and vascular pseudo-diffusion \citep{langkilde2018evaluation}.
Accordingly, a large family of richer models has been proposed to recover this microstructural detail, including intravoxel incoherent motion (IVIM) \citep{lebihan1988separation}, biexponential \citep{mulkern2006biexponential} and multi-exponential \citep{conlin2021improved} fitting, and biophysical models such as VERDICT \citep{panagiotaki2014microstructural}, which varies diffusion time, and hybrid multidimensional MRI \citep{chatterjee2018hybrid}, which adds relaxation contrast.

A consistent pattern has emerged from this literature. The added compartments carry biologically meaningful information, yet they rarely outperform ADC for the detection task itself, with the clearer gains appearing in grade stratification.
\citet{langkilde2018evaluation}, on the extended-b cohort re-analyzed here, found compartment fractions complementary to but not better than ADC for tissue characterization. The VERDICT intracellular volume fraction outperformed ADC specifically in separating lesions containing Gleason pattern 4 from benign and Gleason 3+3 tissue \citep{johnston2019verdict}. A recent head-to-head comparison of monoexponential, fractional-order-calculus, and multi-compartment models reported comparable performance for distinguishing cancer from benign hyperplasia while the multi-compartment fractions improved risk stratification \citep{he2025improved}.
The task and the comparator matter. For patient-level, delineation-free detection, compartment-isolating techniques such as restriction spectrum imaging do outperform whole-gland minimum ADC \citep{domingo2025restriction}, whereas the comparisons above---and the present study---concern localized detection in a delineated region of interest (ROI).
Within that setting the recurring near-tie invites a different question. Rather than proposing another compartment model in the hope of beating ADC, we ask why a single scalar is so hard to beat, and where added contrast, rather than added model complexity, would have to come from.

A fundamental challenge in multi-compartment diffusivity estimation is the ill-posedness of the inverse Laplace transform. Recovering a spectrum of diffusivities from a finite, noisy set of b-value measurements is inherently underdetermined \citep{prange2009quantifying}.
Point estimation methods, including non-negative least squares and regularized regression, yield a single spectrum without quantifying how reliably each component is determined by the data. Bayesian approaches address this by sampling the full posterior, providing per-component uncertainty estimates that distinguish well-constrained spectral features from those effectively unconstrained by the measurements \citep{prange2009quantifying}.

Building on the diffusivity-spectrum estimator of \citet{wells2022estimation}, we present a Bayesian spectral decomposition framework for multi-b prostate DWI, apply it to 56 patients (149 ROIs), and estimate spectra by both a maximum a~posteriori (MAP) point estimate and full Bayesian inference via the No-U-Turn Sampler (NUTS) \citep{hoffman2014nuts}.
We investigate three questions. First, how well each diffusivity component can be measured, both theoretically (Fisher information and the Cram\'er-Rao bound) and practically (the NUTS posterior). Second, how spectral classifiers compare with ADC for ROI-level tumor detection. Third, whether their near-equivalence can be explained by an alignment between ADC's spectral sensitivity and the optimal classification direction. Grade-associated spectral changes are examined descriptively, the graded subgroup being too small for a performance comparison.
We additionally demonstrate the framework at the pixel level in a single patient from a separate, earlier acquisition, producing spectral fraction maps, a discriminant score, and per-voxel uncertainty estimates.
